\bigskip
\begin{center}\bfseries
III. Die historische Jugendperiode: 1796 bis zur Doktorpromotion am 16. Juli 1799.
\end{center}
\begin{center}\textls[120]{a) Die ersten Studiensemester in Göttingen.}\end{center}
\indent Wie wir gesehen haben, besitzen wir über den Stand von \textsc{Gauss}’
analytischen Studien vor dem Abgang an die Universität keine weiteren Anhaltspunkte
als seine Angabe des Jahres 1791 für das agM., ferner für 1794 die
mündliche Überlieferung über den Zusammenhang zwischen agM. und den
Potenzreihen, deren Exponenten die Quadratzahlen sind, und die Tatsache der
Erwerbung der \textit{Prinzipien} \textsc{Newtons}. Dass diese in Braunschweig das einzige
klassische Werk waren, das \textsc{Gauss} zur Verfügung stand, ist sehr wahrscheinlich;
umso grösser mag seine Begierde nach literarischen Hilfsmitteln gewesen
sein, und da Göttingen diese in höherem Mass zu bieten vermochte, als die
braunschweigische Landesuniversität Helmstedt, entschied sich \textsc{Gauss} für die
hannöversche Hochschule. Am 21. August 1795 erging eine Herzogliche Verfügung,
dass dem nach Göttingen abgehenden Studiosus \textsc{Gauss} während seines
Studiums daselbst eine Unterstützung von 158 Talern jährlich und ein Freitisch
in Göttingen zugestanden sei; am 11. Oktober 1795 reiste \textsc{Gauss} nach
Göttingen ab und wird daselbst am 15. Oktober als »matheseos cult.« immatrikuliert.
— Im ersten Semester schwankt \textsc{Gauss} noch, ob er sich den alten
Sprachen oder der Mathematik und Astronomie widmen solle. Mathematik
lehrte damals \textsc{Abraham Gotthelf Kaestner} (1719—1800), »der erste Mathematiker
unter den Dichtern und der erste Dichter unter den Mathematikern«,
wie \textsc{Gauss} einst scherzhaft von ihm sagte, Astronomie vertrat \textsc{Carl Felix
Seyffer} (1762—1824), mit dem \textsc{Gauss} später in freundschaftlichem Verkehr
gestanden hat, Physik, \textsc{G. Christian Lichtenberg} (1744—1799), klassische
Philologie, \textsc{Chr. Gottlob Heyne} (1729—1812) und Geschichte, \textsc{Arnold Heeren}
(1760—1842).

\clearpage\SourceLeaf{104}\ArithmeticHead{19}{R}{III a. DIE ERSTEN STUDIENSEMESTER IN GÖTTINGEN.}
Über seine ersten Eindrücke in Göttingen berichtet \textsc{Gauss} wenige Tage
nach seiner Ankunft in einem Briefe an seinen Braunschweiger Lehrer und
Gönner \textsc{Zimmermann}\footnote{Mit noch zwei spätern Briefen an denselben Empfänger veröffentlicht in der Braunschweiger G. N. C. Monatsschrift, 1921, IX.} vom 19. Oktober mit den folgenden Worten:

»Die Aufnahme, die ich zum Teil hier gefunden habe, könnte mich unruhig
machen, wenn ich nicht darauf vorbereitet gewesen wäre. Von \textsc{Heyne}
glaube ich gut aufgenommen zu sein; es scheint, als wenn er sich für mich
interessiert. Wegen des philologischen Seminarium habe ich bisher noch
nicht mit ihm sprechen können. In \textsc{Kaestner} glaubte ich anfangs einen
stumpfen Greis zu finden, von dem ich mir keine tätige Unterstützung versprechen
könnte. Allein schon jetzt sehe ich, dass ich mich geirrt habe, und
dass er ein sehr vortrefflicher Mann ist. In Prof. \textsc{Heeren} habe ich einen
ungemein liebenswürdigen und gefälligen Mann kennen gelernt. Prof. \textsc{Seyffer}
ist nach Schwaben verreist; auch der Hofr. \textsc{Lichtenberg} hält sich nicht in
Göttingen auf.

Ich habe die Bibliothek besehen und ich verspreche mir davon einen
nicht geringen Beitrag zu meiner glücklichen Existenz in Göttingen. Ich
habe schon mehrere Bände von den Comment[ariis] Acad[emiae Scientiarum]
Petrop[olitanae] im Hause und eine noch grössere Anzahl habe ich durchblättert.
Ich kann nicht leugnen, dass es mir sehr unangenehm ist zu finden,
dass ich den grössten Teil meiner schönen Entdeckungen in der unbestimmten
Analytik nur zum zweiten Male gemacht habe. Was mich tröstet ist dieses.
Alle Entdeckungen \textsc{Eulers}, die ich bis jetzt gefunden habe, habe ich auch
gemacht, und noch einige mehr. Ich habe einen allgemeinern und wie ich
glaube natürlichern Gesichtspunkt getroffen; ich sehe noch ein unermessliches
Feld vor mir und \textsc{Euler} hat seine Entdeckungen in einem Zeitraume von
vielen Jahren nach manchen vorher gegangenen \textls[120]{tentaminibus} gemacht.

Den Freitisch habe ich seit einigen Tagen gehabt: das Essen ist leidlich
.\ .\ .\ .\ .\kern.12em.«

Der ganz unter \textsc{Kaestners} Einfluss stehende Betrieb der mathematischen
Vorlesungen wird auf \textsc{Gauss}’ Entschluss, zwischen der Philologie und der
Mathematik zu entscheiden, kaum eingewirkt haben, vielmehr soll\footnote{Siehe \textsc{Sartorius v. Waltershausen}, a. a. O.\textsuperscript{13)}, S. 16.} die Ent-\par
\HistoricalFoot{}{3*}
\clearpage\SourceLeaf{105}\ArithmeticHead{20}{L}{L. SCHLESINGER, ÜBER GAUSS’ FUNKTIONENTHEORETISCHE ARBEITEN.}
\noindent scheidung erst am 30. März 1796 in Braunschweig gefallen sein, an dem Tage,
wo \textsc{Gauss} seine erste bedeutsame mathematische Entdeckung gemacht hat,
indem er als Folge einer allgemeinen Theorie der Kreisteilung erkannt hat,
dass der Kreisumfang mit Zirkel und Lineal in 17 gleiche Teile zerlegt werden
kann. Mit diesem Tage und mit dieser Entdeckung beginnt \textsc{Gauss} auch sein
wissenschaftliches \textit{Tagebuch}\footnote{Werke X, 1, S. 488; ebenda S. 125, in einem Briefe an \textsc{Gerling} vom 6. Jan. 1819, gibt \textsc{Gauss} selbst einen Bericht über das Geschichtliche seiner Entdeckung.} und damit das erste Blatt zur Geschichte seiner
eigenen wissenschaftlichen Entwicklung. Auf Veranlassung \textsc{Zimmermanns} veröffentlichte
\textsc{Gauss} seine Entdeckung im Intelligenzblatt der allgemeinen Literaturzeitung
(Werke X, 1, S. 3) und machte nach seiner Rückkehr nach
Göttingen (die nach dem \textit{Tagebuch} zwischen dem 12. und 29. April 1796 erfolgte)
auch \textsc{Kaestner} davon Mitteilung; wie wenig Verständnis er aber bei
diesem fand, zeigt der folgende Brief an \textsc{Zimmermann} vom 26. Mai 1796:

»Ich müsste mich sehr irren (und das scheint mir doch bei einem so
freien offnen Manne schwer) wenn die Kälte, mit der der H. H[ofrath] K[\textsc{aestner}]
meine Entdeckung aufnahm, aus irgend einer persönlichen Rücksicht auf mich,
und nicht vielmehr aus einer gewissen Abneigung gegen alles Neue, die man
auch bei der grössten Hochachtung an ihm nicht verkennen kann, herrührte.
Seinem nachherigen Benehmen nach scheinen auch die Umstände, in denen
er damals war, Anteil daran gehabt zu haben. Als ich einige Zeit nachher
wieder bei ihm war, fing er selbst von dem Satze an, behauptete aber ganz
positiv, dass die Natur der Gleichungen die Unmöglichkeit der Sache beweise.
Ich stellte ihm wieder (wie das erste Mal) vor, dass die Gleichung nicht wie
er glaubte vom 17ten sondern, weil man eine Wurzel sogleich von selbst kenne,
nur vom 16ten Grade sei, und dass diese sich auf eine vom 8ten, diese auf
eine vom 4ten, 2ten [reduzieren] u. also auflösen lasse. Das schien ihn dieses
Mal zu frappieren und so nahm er sogleich einen andern Ton an und sagte,
wenn es auch richtig wäre, so würde es doch gar keinen Nutzen haben, da
man die Vielecke weit leichter nach den Tafeln konstruieren könne, und endlich
wandte er es so, dass er behauptete, der Grund von der ganzen Sache
sei schon in seinen \textit{Anfangsgründen}\footnote{A. G. \textsc{Kaestner}, \textit{Anfangsgründe der Analysis endlicher Grössen}. Göttingen, Vandenhok 1794,} und er habe es nur nicht der Mühe
wert gehalten, es zu entwickeln. Das scheint mir doch inkonsequent. Da ich\par
\clearpage\SourceLeaf{106}\ArithmeticHead{21}{R}{III a. DIE ERSTEN STUDIENSEMESTER IN GÖTTINGEN.}
\noindent das erste zugab und in dem andern nicht widersprach, so erhielt ich endlich
die Erlaubnis, es ihm schriftlich vorzulegen.

Die Sache selbst verhält sich so. Aus meinem Verfahren kann man zwar mit
vieler Mühe eine Auflösung der Gleichung herleiten, allein diese Auflösung
directe und wie er jetzt behaupten wollte, ist zuverlässig unmöglich. Das
wird er auch aus dem Aufsatze, den ich ihm vorlegte, einigermassen gesehen
haben. Ausser einigen unbedeutenden Erinnerungen, die das Wesentliche
der Sache nicht betreffen, fiel dann sein Urteil so aus, »dass sie praktischen
Nutzen keinen haben würde, dass sie aber, zumal, wenn ich die Sache noch
allgemeiner behandeln könne, eine ganz artige Kuriosität sei und vielleicht
eine hellere Einsicht in diesen Teil der Mathematik verschaffen werde«, welches
Urteil ich ganz beherzige. Er setzte noch hinzu, dass es doch nicht
ganz neu sei, da er schon gelehrt habe, dass es auf die Auflösung der ofterwähnten
Gleichung ankäme, welchen Glauben ich ihm gern gönnen will.

Ich muss noch etwas über den Fortgang meiner eigentlich analytischen
Untersuchungen sagen. Ich habe in den Mémoires von Paris 1785 eine vortreffliche
Abhandlung von \textsc{Le Gendre} gefunden, wo ein Beweis von dem Lehrsatze,
den ich so lange \textls[120]{vollständig} zu beweisen umsonst gesucht hatte,
vorkommt, aber, wo gerade das angenommen wird (je ne suppose que ce que
.\ .\ .\ sagt L. G.), was allein mir seit beinahe einem Jahre noch fehlte und
was ich nunmehr gefunden habe. Auch habe ich von \textsc{La Granges} Zusätzen
zu \textsc{Eulers} Algebra (worauf Sie mich schon aufmerksam machten, die ich hier
aber nicht habe bekommen können und die auch H. H. \textsc{Kaestner} nicht kannte)
jetzt durch Herrn H. \textsc{Kaestners} Güte eine deutsche Übersetzung in Händen,
die in diesem Jahre erst herausgekommen ist und einen Hrn. Hofrat \textsc{Kausler}
in Württemberg zum Verfasser hat. Dieses an sich vortreffliche Buch enthält
indes wenig, was mit meinen Entdeckungen in Kollision käme, den oft erwähnten
Lehrsatz hat \textsc{La Grange} nicht einmal recht gekannt, vielweniger
bewiesen und \textsc{Le Gendres} Abhandlung ist angeführt, aber ihr Verdienst ganz
schief beurteilt. Es wäre mir sehr daran gelegen, das Jahr zu wissen, wo das
französische Original herausgekommen ist. Der Übersetzer hat auch einige
Zusätze gemacht, die aber von geringem Gehalte sind. Ich glaube es also
jetzt wagen zu können, wenn ein Buchhändler sich findet, an die Ausarbeitung
zu gehen. .\ .\ .\ .«

\clearpage\SourceLeaf{107}\ArithmeticHead{22}{L}{L. SCHLESINGER, ÜBER GAUSS’ FUNKTIONENTHEORETISCHE ARBEITEN.}
Der Lehrsatz, von dem im letzten Absatz dieses Briefes die Rede war,
ist das Reziprozitätsgesetz der quadratischen Reste, dessen strenger Beweis
\textsc{Gauss} am 8. April in Braunschweig geglückt ist, wie die zweite Eintragung
in das \textit{Tagebuch} (Werke X, 1, S. 489, vergl. \textsc{Bachmann}, Abh. 1 dieses Bandes
S. 5) bezeugt. Wir sehen diesen merkwürdigen Studenten nicht nur von Entdeckung
zu Entdeckung fortschreiten, sondern auch eifrig bemüht, sich mit
der Literatur seiner Wissenschaft bekannt zu machen, und wir finden vom
Frühjahr 1796 ab zahlreiche Spuren seiner Beschäftigung mit den Schriften
nicht nur von \textsc{Euler} und \textsc{Lagrange}, sondern auch mit denen der \textsc{Bernoullis},
\textsc{Stirlings} und \textsc{Vandermondes}. Das Verzeichnis \textsc{Euler}scher Schriften auf den
ersten Durchschussblättern des \textit{Leiste} wurde schon erwähnt, im Nachlass finden
sich aus dieser Zeit stammende Auszüge aus den Schriften der Akademien
von Paris, Berlin, St. Petersburg, Turin, auf Buchdeckeln, in Notizheften und
auf Zetteln stehen Lesefrüchte, besonders aus \textsc{Euler} und \textsc{Lagrange}. — Ein
Bild von \textsc{Gauss} aus diesen Studentenjahren ist uns in Briefen und Aufzeichnungen
des siebenbürgischen Mathematikers \textsc{Wolfgang Bolyai von Bolya}
(1775—1856) erhalten, der von Oktober 1796 bis Juni 1799 in Göttingen
studiert, und der \textsc{Gauss} im Hause des Professors \textsc{Seyffer} kennen gelernt hat.
In einer 1840 verfassten Selbstbiographie\footnote{Original magyarisch, vergl. \textit{Briefwechsel Gauss-Bolyai}, S. 178.} berichtet \textsc{Bolyai}:

»Wir [nämlich er und ein Landsmann, \textsc{Baron Kemény}] gingen nach Göttingen,
wo .\ .\ .\ .\ ich mit dem damals dort studierenden \textsc{Gauss} bekannt wurde,
mit dem ich noch heute in Freundschaft bin .\ .\ .\ .\quad Er war sehr bescheiden
und zeigte wenig; nicht drei Tage wie mit \textsc{Plato}, Jahre lang konnte man
mit ihm zusammen sein, ohne seine Grösse zu erkennen. Schade, dass ich
dieses titellose, schweigsame Buch nicht aufzumachen und zu lesen verstand!
.\ .\ .\ .\ Uns verband die sich äusserlich nicht zeigende Leidenschaft für Mathematik
und unsere sittliche Übereinstimmung«. Und in einem nach \textsc{Gauss}’ Tode
an den Biographen von \textsc{Gauss}, den Göttinger Professor \textsc{Sartorius von Waltershausen}
gerichteten Briefe \textsc{Bolyais}\footnote{A. a. O., S. 152.} heisst es: »Von Jena ging ich nach Göttingen,
wo ich den sehr gütigen Professor \textsc{Seyffer} besuchend \textsc{Gauss} zuerst
sah; ich .\ .\ sprach dreist .\ .\ über die Seichtigkeit der Behandlung der Gründe
der Mathematik in Hinsicht der Multiplikation, Division, Potenz, geraden\par
\clearpage\SourceLeaf{108}\ArithmeticHead{23}{R}{III a. DIE ERSTEN STUDIENSEMESTER IN GÖTTINGEN.}
\noindent Linie, Ebene, Gleichheit in verschiedener Hinsicht u. dgl. Nach diesem begegneten
wir uns auf dem Walle, jeder war allein, gesellten uns, gingen zu
einander und bald schwuren wir unter der Fahne der Wahrheit Brüderschaft.
Hierauf ruhte er von seiner anhaltenden stillen Arbeit meistens bei mir aus;
sprach nie im Voraus, selbst bei fertigem schweigend.« An seinen Sohn \textsc{Johann}\footnote{Original magyarisch, vergl. \textsc{Stäckel}, \textit{Die beiden Bolyai} I, 1913, S. 90. Dieses Werk orientiert auch am besten über die Beziehungen von \textsc{Gauss} zu \textsc{Bolyai} Vater und Sohn. Siehe auch \textsc{Stäckels} Abhandl. 4 in diesem Bande, Abschnitt I, S. 6—46.}
schreibt \textsc{Bolyai} um 1825: ».\ .\ Ich kann Dir es schriftlich zeigen, dass
er [\textsc{Gauss}] sich über die Parallelen den Kopf zerbrochen hat. .\ .\ Meine Ideen
haben ihm sehr gut gefallen, obwohl sie ihn bei weitem nicht befriedigten,
und er machte mich darauf aufmerksam, eine wie wichtige Materie die Sache
der Parallelen sei. In den Elementen der Arithmetik und Geometrie war
\textsc{Gauss} damals (obwohl sonst viele Turmstockwerke über mich hinausragend)
weniger stark als ich, blos durch mich selbst, aber ihm waren die höheren
Rechnungen nur noch ein Spiel, wo ich noch nicht einmal ein Ahnung davon
hatte« .\ .\ .\ Dass \textsc{Gauss} in diesen »höheren Rechnungen« auch schon zu
selbständigen Forschungen vorgedrungen war, wenn auch zu jener Zeit die
Zahlentheorie noch vorwiegend seine Gedankenwelt beherrscht haben dürfte,
zeigt das \textit{Tagebuch} und andere Aufzeichnungen, die \textsc{Gauss} aufbewahrt hat.

Die Aufzeichnung 7. des \textit{Tagebuchs} vom 24. Mai 1796, die erste, die sich
auf einen ausgesprochen analytischen Gegenstand bezieht (Umwandlung einer
unendlichen Reihe in einen Kettenbruch) steht einerseits, wie die Aufzeichnung
58. zeigt\footnote{Die Notiz 58. vom 16. Febr. 1797 knüpft ausdrücklich an die Notiz 7. an; sie gibt die Umwandlung der Reihe
\[
1-a+a^3-a^6+a^{10}-\ldots
\]
in einen Kettenbruch und schliesst mit den Worten: »wodurch leicht alle Reihen, deren Exponenten eine Reihe zweiter Ordnung bilden, transformiert werden«. Es ist die älteste uns erhaltene Aufzeichnung, in der diese Art von Reihen auftritt.},
mit den Potenzreihen im Zusammenhang, deren Exponenten
eine arithmetische Progression zweiter Ordnung bilden, andererseits knüpft
sie, wie eine Eintragung in \textsc{Lamberts} \textit{Tafeln} (abgedruckt Werke X, 1, S. 491)
erkennen lässt, unmittelbar an \textsc{Eulers} Abhandlung \textit{de seriebus divergentibus}\footnote{Novi Comm. Acad. Petropol. 5 (1754/5), 1760, S. 205.}
an. Dass es sich hier nicht um ein nur vorübergehendes Interesse handelt,\par
\clearpage\SourceLeaf{109}\ArithmeticHead{24}{L}{L. SCHLESINGER, ÜBER GAUSS’ FUNKTIONENTHEORETISCHE ARBEITEN.}
\noindent geht daraus hervor, dass \textsc{Gauss} auf einen besondern Fall der hier vorkommenden,
von \textsc{Euler} betrachteten Reihe \(1+\displaystyle\sum_{k=0}^{\infty}(-1)^{k+1}m(m+n)\ldots(m+kn)x^{k+1}\) Ende
des Jahres 1797 in einer ausführlichen Untersuchung (siehe Werke X, 1,
S. 382) zurückgekommen ist, während er die Umformung der allgemeinen
\textsc{Euler}schen Reihe in einen Kettenbruch in dem Art. 14 der Abhandlung
\textit{circa seriem} von 1812 (Werke III, S. 138, Gl. (36)) als Beispiel angibt.

In den Tagebuchnotizen 9., 13., 14., 31. (31. Mai bis 6. September 1796)
behandelt \textsc{Gauss} sogenannte mittlere Werte der Zahlentheorie; in einer aus viel
späterer Zeit (Ende 1799) stammenden Aufzeichnung in der Scheda Ac (abgedruckt
Werke X, 1, S. 14—16) wird neben solchen zahlentheoretischen
Mittelbildungen ein Grenzwert betrachtet, der als eine Art von harmonisch-arithmetischem
Mittel bezeichnet werden kann, und auf den in der Handschrift
unmittelbar die Notiz (Werke X, 1, S. 551 oben) über das harmonisch-geometrische
Mittel folgt, das sich auf das agM. zurückführen lässt. Es zeigt
dies, dass für \textsc{Gauss}, wenigstens zu jener Zeit, alle diese Mittelbildungen aus
einem gemeinsamen Interesse herauswuchsen, andererseits zeigen diese Untersuchungen
das erwachende Bestreben, asymptotische Werte zu betrachten, das, wie
wir sehen werden, noch weiterhin andauert.

Aus dem August 1796 stammt ferner das \textit{Exercitationes Mathematicae} überschriebene
Blättchen (Werke X, 1, S. 138, vergl. auch \textit{Tagebuch} 28., ebenda,
S. 500), das durch die Vielseitigkeit seines Inhalts ein Bild von der Vielseitigkeit
der Interessen des jungen \textsc{Gauss} gibt. Wenn \textsc{Gauss} im Art. 3 bei einer
Verallgemeinerung der \textsc{Leibniz}schen Reihe \(1-1+1-\ldots\) die »metaphysische
Schlussweise« \textsc{Daniel Bernoullis} unbefriedigend findet und sie durch eine mathematische
ersetzt, so sehen wir, dass er auch in der Analysis schon den kritischen
Geist walten lässt, den er in der Zahlentheorie beim Reziprozitätsgesetz, in
der Algebra beim Beweise der Wurzelexistenz, in der Geometrie beim Parallelenaxiom
zur Geltung bringt. Wenn er in den Artikeln 5, 6 für die Integrale
\[
\int\limits_0^x\frac{dx}{\sqrt{1-x^3}}\mathpunct{\vcenter{\hbox{,}}}
\qquad
\int\limits_0^x\frac{dx}{\sqrt{1-x^4}}\mathpunct{\vcenter{\hbox{,}}}
\]
als den nächsten Fällen nach dem arc sin, die obere Grenze als Funktion
des Integrals in Reihen entwickelt, so zeigt das, dass er auch in der Integral-\par
\clearpage\SourceLeaf{110}\ArithmeticHead{25}{R}{III a. DIE ERSTEN STUDIENSEMESTER IN GÖTTINGEN.}
\noindent rechnung auf der Suche nach neuen Aufgaben mit der unbewussten Sicherheit
des Genius die Goldader zu finden weiss. Freilich arbeitet er hier, wie auch
in den Tagebuchnotizen 32. und 33. vom 9. und 14. September (Werke X, 1,
S. 502), noch mit den ganz formalen Methoden der Reihenentwicklung und
Inversion, die er sich wohl noch in Braunschweig aus \textsc{Newtons} \textit{Prinzipien}
angeeignet und später durch das Studium von \textsc{Stirlings} \textit{Methodus Differentialis}
befestigt hat\footnote{In bezug auf die bei Art. 8 der \textit{Exercitationes} auftretende Figur (sich überschlagendes Fünfeck, sog. Drudenfuss) bemerken wir, dass dieselbe, wie \textsc{Gauss} 1825 in einem Briefe an \textsc{Olbers} andeutet (Werke VIII, S. 398), bei \textsc{Meister} vorkommt; vergl. Werke XII, S. 53, 54.}.
Dieselben Methoden, gewissermassen ein formales Rechnen
mit Potenzreihen, liegen den Tagebuchnotizen 29. (Sinus und Kosinus höherer
Ordnung) und 45. (Umwandlung von \(\displaystyle\sum_{n=1}^{\infty}\frac{(-1)^{n-1}}{n^\omega}\) in eine Potenzreihe von \(\omega\)) zu
Grunde, die andererseits wiederum auf das Studium \textsc{Euler}scher Schriften hinweisen,
und auch die Tagebuchnotiz 49., in der ein Beweis der \textsc{Lagrange}schen
Umkehrungsformel angekündigt wird, gehört in diesen Gedankenkreis.
Der Satz von \textsc{Lagrange} hat auf \textsc{Gauss} eine ganze besondere Anziehung ausgeübt,
wohl weil er ihm ein brauchbares Werkzeug für das Rechnen mit
Reihen lieferte. \textsc{Gauss} ist noch wiederholt auf diesen Satz zurückgekommen,
so im \textit{Tagebuch} Nr. 86 vom Mai 1798, in Aufzeichnungen bei S. 5, 10—12
und 19 des \textit{Leiste}, und in den Werke VIII, S. 76—83 aus dem Nachlass
abgedruckten Notizen. — Die ältere dieser Notizen (S. 76—79) ist dadurch
besonders bemerkenswert, dass sie bestimmt war, die erste grössere Veröffentlichung
von \textsc{Gauss} zu werden. \textsc{Kaestner} hat sie etwa im April 1797 an
\textsc{Hindenburg}, den Herausgeber des Archivs der reinen und angewandten Mathematik,
zur Publikation gesandt, aber die Arbeit des unbekannten Anfängers
blieb liegen und ist erst 1899 im Bande VIII der Werke an die Öffentlichkeit
gelangt (vergl. Werke X, 1, S. 429, 443, 444, 509). — Die spätere Notiz
(S. 80—83) stammt aus dem Handbuch 19, Be, S. 273—276 und trägt das Datum
13. Mai 1847, ist aber wie die älteren Versuche und wie die an \textsc{Hindenburg}
gesandte Schrift rein formal; zu einer Konvergenzuntersuchung findet sich kein
Ansatz. Dagegen zeigen die Zitate in dem zur Veröffentlichung bestimmten
Aufsatze (Werke VIII, S. 76) vollständige Beherrschung der einschlägigen
Literatur, namentlich wird die Abhandlung von \textsc{Lagrange} aus dem 24. Bande\par
\HistoricalFoot{X\textsubscript{2} Abh. 2.}{4}
\clearpage\SourceLeaf{111}\ArithmeticHead{26}{L}{L. SCHLESINGER, ÜBER GAUSS’ FUNKTIONENTHEORETISCHE ARBEITEN.}
\noindent der Berliner Memoiren von 1768—1770 angeführt. — Man kann daraus
schliessen, dass \textsc{Gauss} zu jener Zeit auch die im 25. Bande dieser Memoiren
enthaltene Arbeit von \textsc{Lagrange} \textit{Sur le Problème de Képler} studiert haben
wird, die ja eine der bedeutsamsten Anwendungen des Umkehrungssatzes enthält.
Die Spuren der Beschäftigung mit dieser Arbeit lassen sich in späteren
Untersuchungen von \textsc{Gauss} deutlich nachweisen. — \textsc{Lagrange} entwickelt die
exzentrische und die wahre Anomalie, sowie den Radiusvektor nach Sinus
und Kosinus der Vielfachen der mittleren Anomalie, und drückt die Koeffizienten
dieser Entwicklung durch Reihen aus, die nach Potenzen der Exzentrizität
fortschreiten. \textsc{Gauss} hat sich etwa 1805 mit dem asymptotischen Verhalten
dieser Koeffizienten eingehend beschäftigt (Werke X, 1, S. 420ff., vergl.
auch die Bemerkungen S. 444), wohl im Verfolg von Untersuchungen über
asymptotische Gesetze der Zahlentheorie, wie sie die bereits genannten Tagebuchnotizen
aus der Zeit vom 31. Mai bis 6. September 1796 anzeigen. —
Bekanntlich wird das infinitäre Verhalten der Koeffizienten in den Entwicklungen
der exzentrischen Anomalie und des Radiusvektors durch die \textsc{Bessel}sche
Funktion \(J_n(ne)\) und ihre Derivierten für grosse Werte von \(n\) dargestellt;
die Notiz 82. des \textit{Tagebuchs} (Werke X, 1, S. 385, 388) zeigt, dass \textsc{Gauss} im
Oktober 1797 mit dem asymptotischen Verhalten der Wurzeln der Gleichung
\(J_1(x)=0\) beschäftigt war. — Eine weitere asymptotische Wertbestimmung
enthält die Tagebuchnotiz 83. vom April 1798 (Werke X, 1, S. 525), bei der
die Bedeutung der dort auftretenden Funktion \(l(1+x)\), aus dem von \textsc{Gauss}
angegebenen asymptotischen Ausdruck ihrer \(i\)ten Iterierten für grosse Werte
von \(i\) nicht vollständig bestimmt werden konnte. Auch bei diesen Untersuchungen
hat \textsc{Gauss} noch mit formal angesetzten Reihen operiert, durch die
er Funktionalgleichungen zu befriedigen suchte; es ist anzunehmen, dass er
seine zahlentheoretischen Grenzausdrücke auch durch ähnliche Methoden heuristischer
Art ermittelt haben wird. —

Es ist jetzt noch eine weitere Anregung zu besprechen, die \textsc{Gauss} aus
Arbeiten von \textsc{Lagrange} geschöpft hat. Wenn \(a,b\) die grossen Achsen der
Bahnen zweier Planeten bedeuten, so treten in der Entwicklung der Störungsfunktion
Ausdrücke von der Form auf:
\[
(a^2+b^2-2ab\cos\theta)^{-n},\quad n=\tfrac12,\tfrac32,\tfrac52\ldots,
\]
\clearpage\SourceLeaf{112}\ArithmeticHead{27}{R}{III a. DIE ERSTEN STUDIENSEMESTER IN GÖTTINGEN.}
\noindent die ihrerseits wieder nach den Kosinus der Vielfachen von \(\theta\) zu entwickeln
sind. Die Koeffizienten dieser Entwicklungen können selbst wieder nach Potenzen
von \(b/a=t\) entwickelt werden, wenn \(b<a\) ist, es sind die sogenannten
Funktionen der grossen Achsen oder Koeffizienten von \textsc{Laplace}. — Diese
Art von Entwicklungen findet sich in Arbeiten von \textsc{Euler} und \textsc{Lagrange}\footnote{L. \textsc{Euler}, \textit{Pièce qui a remporté le prix etc. en 1748} (Paris 1749) und \textit{Inst. calc. integralis} I (1768) § 279, Opera, ser. I, vol. 11, S. 165; \textsc{Lagrange}, Misc. Taurin. 1762/65, Oeuvres I, S. 620, Recueil des pièces qui ont remporté les prix etc. t. 9 1766, Oeuvres I, S. 88; vergl. die Bemerkung Werke X, 1, S. 531.};
der letztere zerlegt
\[
(1-2t\cos\theta+t^2)=(1-te^{i\theta})(1-te^{-i\theta})
\]
und entwickelt dann die beiden Binome. Dieses Verfahren reproduziert \textsc{Gauss}
im Art. 6 der Abhandlung \textit{circa seriem} (1812, Werke III, S. 128), was aus
dem Grunde schon hier hervorgehoben werden muss, weil \textsc{Gauss} die Gewohnheit
hat, in spätern Arbeiten ältere Methoden und Resultate, oft auch ohne
direkten Zusammenhang mit dem daselbst behandelten Gegenstande, wiederzugeben
und auf diese Weise gleichsam für sich festzulegen. Nun sind die
Koeffizienten der Entwicklung von \((1-2t\cos\theta+t^2)^{-n}\) nach den Kosinus der
Vielfachen von \(\theta\) als Funktionen von \(t^2\) durch die sogenannte \textsc{Gauss}sche Reihe
\(F(\alpha,\beta,\gamma,x)\) darstellbar und lassen sich insbesondere für \(n=\frac12\) durch Integrale
darstellen, die, wie wohl zuerst \textsc{d’Alembert} (1754) hervorgehoben hat, sich
»durch die Rektifikation der Kegelschnitte integrieren lassen«, endlich liefert
für ein beliebiges \(n\) die Transformation, die \textsc{Gauss} a. a. O. (Werke III, S. 129)
unter »tertio fit« angibt und die darin besteht, dass die nach Potenzen von
\(t^2=b^2/a^2\) fortschreitenden Reihen in solche nach Potenzen von \(4ab/(a+b)^2\)
umgeformt werden, direkt den Algorithmus des agM. — Wir erkennen also
in dieser Entwicklung der Störungstheorie gewissermassen ein Sammelbecken
für die wesentlichsten analytischen Probleme, die \textsc{Gauss} sein ganzes Leben
hindurch begleitet haben, nämlich
\begin{list}{}{\setlength{\leftmargin}{5em}\setlength{\itemsep}{0pt}\setlength{\parsep}{0pt}}
\item[1.] die Reihe \(F(\alpha,\beta,\gamma,x)\),
\item[2.] die elliptischen Funktionen und Integrale,
\item[3.] das agM.
\end{list}
\noindent Gestützt wird diese Auffassung durch den Umstand, dass gerade die gedachte\par
\HistoricalFoot{}{4*}
\clearpage\SourceLeaf{113}\ArithmeticHead{28}{L}{L. SCHLESINGER, ÜBER GAUSS’ FUNKTIONENTHEORETISCHE ARBEITEN.}
\noindent Entwicklung bezw. Transformation im Falle \(n=\frac12\) die Reihe für \(4/\pi\) liefert,
die \textsc{Gauss} in der Nr. 87 des \textit{Tagebuchs} (Juni 1798) und in seinem Exemplar
von \textsc{Schulzes} \textit{Tafeln} aufgezeichnet hat (vergl. die bereits genannte Bemerkung
Werke X, 1, S. 531, 532), ferner ist der Zusammenhang, der in diesem Falle
zwischen den Koeffizienten und den elliptischen Integralen besteht, in der aus
dem Jahre 1800 stammenden nachgelassenen Abhandlung über das agM. (Werke
III, S. 370, 371) ausführlich behandelt und auf die Bedeutung für die physische
Astronomie und die Theorie der Planetenstörungen ausdrücklich hingewiesen.
— Auch in dem oben erwähnten Briefe an \textsc{Schumacher} (Werke X, 1,
S. 247) wird gerade dieser Teil der Abhandlung von 1800 besonders hervorgehoben. —

Nachdem wir so den Einflüssen nachgegangen sind, die die Lektüre der
Schriften von \textsc{Euler} und \textsc{Lagrange} auf die Ideenentwicklung von \textsc{Gauss} ausgeübt
haben dürfte, wenden wir uns jetzt der Aufgabe zu, diese Ideenentwicklung
selbst in den Göttinger Studentenjahren zu schildern.

Zu Anfang des Jahres 1797 beginnt \textsc{Gauss} sich systematisch mit der
Theorie der elliptischen, d. h. zunächst der lemniskatischen Funktionen zu
beschäftigen. Dieser Gegenstand hat ihn in Verbindung mit der Untersuchung
der allgemeineren Transzendenten, die aus der \textsc{Gauss}schen Reihe
\[
F(\alpha,\beta,\gamma,x)=1+\frac{\alpha\,.\,\beta}{1\,.\,\gamma}x
+\frac{\alpha(\alpha+1)\beta(\beta+1)}{1\,.\,2\,.\,\gamma(\gamma+1)}x^2+\ldots
\]
entspringen, viele Jahrzehnte hindurch, wenn auch mit Unterbrechungen gefesselt.
Obwohl \textsc{Gauss} über diese seine Untersuchungen nur verhältnismässig
wenig veröffentlicht hat, haben dieselben doch für die Entwicklung seiner
funktionentheoretischen Gedanken eine so ausschlaggebende Rolle gespielt,
dass wir diese Arbeiten mit besonderer Sorgfalt verfolgen müssen.

Ein flüchtiger Blick wird auf den Stand zu werfen sein, in dem sich die
Theorie der elliptischen Integrale befand, als \textsc{Gauss} ihr seine Bemühungen
zuwandte.

Alle Mathematiker vor \textsc{Euler}, die sich mit Integralen von der Art beschäftigt
haben, die wir heute elliptische nennen, sind von der Aufgabe der
Rektifikation einer krummen Linie ausgegangen. Elastische Kurve und Lemniskate
werden durch das Integral
\clearpage\SourceLeaf{114}\ArithmeticHead{29}{R}{III a. DIE ERSTEN STUDIENSEMESTER IN GÖTTINGEN.}
\[
\int\frac{dx}{\sqrt{1-x^4}}
\]
rektifiziert, Ellipse und Hyperbel durch das Integral
\[
\int\frac{(1-k^2x^2)\,dx}{\sqrt{(1-x^2)(1-k^2x^2)}}\mathpunct{\vcenter{\hbox{.}}}
\]
Das Integral erster Gattung wurde also anfangs nur im lemniskatischen Falle
betrachtet, während im allgemeinen Falle das Integral zweiter Gattung im
Vordergrunde des Interesses stand. Erst \textsc{Euler} brachte in seinen auf das
Additionstheorem bezüglichen Untersuchungen das allgemeine Integral erster
Gattung zur Geltung, wobei zu beachten ist, dass diese Arbeiten \textsc{Eulers}
ihren Ausgang von den noch ganz im geometrischen Bannkreise stehenden
Arbeiten \textsc{Fagnanos} genommen haben, demgemäss mit der Betrachtung des
lemniskatischen Integrals beginnen und sich erst nach und nach von diesem
geometrischen Ausgangspunkt befreien.

Die \textsc{Landen}sche Transformation war zwar von \textsc{Lagrange} in der bereits oben
(Fussnote\textsuperscript{18)}) zitierten Abhandlung auch auf das Integral erster Gattung angewandt
worden, aber im Grunde genommen schien nur die einfache Form, in der
sich bei \textsc{Euler} das Additionstheorem der Integrale erster Gattung dargeboten
hatte, diese Integrale als ein geeignetes Objekt für die weitere Untersuchung
zu bezeichnen.

\textsc{Gauss} hatte, wie wir bemerkt haben, in den Artikeln 4, 5 der \textit{Exercitationes}
(Werke X, 1, S. 140, 141) von vornherein (1796) in gewissen speziellen Integralen
erster Gattung die Quelle für eine Verallgemeinerung der trigonometrischen
Funktionen gesucht; er hat aber damit parallel und durchaus unabhängig
davon das agM. in Verbindung mit den Reihen, deren Exponenten
die Quadratzahlen sind, in den Kreis seiner Betrachtungen gezogen, und wir
werden sehen, wie bei ihm diese beiden Fäden anfangs unabhängig neben
einander herlaufen, bis sie sich — erst verhältnismässig spät, nämlich 1799 —
mit einander vereinigen.

Es möge gleich vorweg ausgesprochen werden, was die im folgenden zu
schildernde Analyse des handschriftlichen Nachlasses ergibt.

\textsc{Gauss} behandelt in den Jahren 1797—1798 das lemniskatische Integral
erster Gattung, in dem er die erste und am nächsten liegende Weiterführung
des trigonometrischen Integrals erkennt, und führt die Theorie der daraus\par
\clearpage\SourceLeaf{115}\ArithmeticHead{30}{L}{L. SCHLESINGER, ÜBER GAUSS’ FUNKTIONENTHEORETISCHE ARBEITEN.}
\noindent entspringenden Funktionen zu einem hohen Grade der Entwicklung; daneben
arbeitet er auch an dem agM. weiter. Der so einfache Zusammenhang zwischen
dem agM. und dem allgemeinen elliptischen Integrale erster Gattung
bleibt ihm lange Zeit verborgen, er entdeckt ihn erst Ende des Jahres 1799
auf einem merkwürdigen Umwege, gewissermassen durch die Beobachtung der
numerischen Übereinstimmung zweier Grössen bis auf die elfte Dezimalstelle.
Von da ab laufen jene beiden Fäden in der Theorie der allgemeinen elliptischen
Funktionen zusammen, in der sie sich gleichsam zu einer höheren Einheit
verbinden. Wir werden der besseren Übersicht wegen diese beiden Fäden,
so lange sie von einander gesondert verlaufen, auch gesondert verfolgen, und
zwar wollen wir zuerst die Theorie der lemniskatischen Funktionen, dann jene
des agM. betrachten.
